% SPDX-License-Identifier: Apache-2.0
\section{Approaches}
\label{sec:approaches}

This section describes each approach and what we found about it.
Section~\ref{sec:comparison} compares them in one table.

\subsection{Userspace USB/IP server in Crostini (A)}

USB/IP sends USB transfers over TCP. A userspace server in \code{penguin}
opens the key with libusb and exports it. The remote host loads
\code{vhci-hcd} and runs \code{usbip attach} through an SSH tunnel. The remote
host then sees the whole key, with all three interfaces.

Two servers can export a real device. jiegec/usbip~\cite{jiegec} (Rust, MIT)
runs one blocking transfer at a time, with a fixed 1\,s timeout, and answers
a cancel request without cancelling anything (verified, source). A YubiKey
report against it, issue~51, closed with no fix (verified).
yunsmall/usbipdcpp~\cite{usbipdcpp} (C++, LGPL-3.0) submits transfers
asynchronously and cancels them correctly (verified, source). It is less than
two years old, and nobody has reported it with a YubiKey (assumption: unknown
maturity). The remote side needs \code{vhci-hcd}, which Debian and Ubuntu
ship as a module~\cite{usbip-deb} (verified).

\subsection{CTAPHID relay with uhid (B)}

A relay moves the 64-byte CTAPHID reports between the real key and a virtual
FIDO key that \code{uhid} creates on the remote host. SSH carries the
reports.

fido-forward~\cite{fidoforward} (Go, MIT) does this. Its key side opens
\code{/dev/hidrawN}. Its remote side, the \term{sink}, runs setuid to a
dedicated user, opens \code{/dev/uhid}, and drops privilege before it reads
any input (verified, source and requirements). Its wire format is a 5-byte
marker, \code{FFWD} and version \code{0x01}, in each direction, then raw
64-byte reports (verified). It filters CTAPHID channels, but not CTAP
commands. The project is one month old and has one author.

qubes-app-u2f~\cite{qubes} (Python, GPL-2.0+) relays parsed CTAP messages,
not raw reports, over Qubes qrexec. It passes only four CTAP2 commands and
sends no keepalives (verified, source). Chromium times out after a PIN
prompt with it, and libfido2 cannot detect it (verified, open issues).

\subsection{Virtual FIDO device with delegated crypto (C)}

A software authenticator on the remote host could send only its signing step
to the real key. This does not work. A real key never exports its private
key, and it signs data that it builds itself: the relying party hash, the
flags, and the counter. bulwarkid/virtual-fido~\cite{virtualfido} signs
inside the library with a Go key (verified, \code{ctap/ctap.go} lines 250 and
364). Delegating to a real key means forwarding whole CTAP commands, which
is approach~B.

\subsection{Agent forwarding (D)}

OpenSSH's agent signs locally, and a forwarded agent lets the remote host ask
for signatures~\cite{openssh} (verified). The FIDO touch happens on the local
key. Since OpenSSH 8.4, a forwarded agent refuses to sign WebAuthn challenges
with FIDO keys~\cite{openssh} (verified), so agent forwarding cannot serve
goal~1. gpg-agent has an \term{extra socket} for forwarding, which refuses
key export and most card commands~\cite{gnupgfwd} (verified). p11-kit can
forward a PKCS\#11 token, such as the YubiKey's PIV applet, over a Unix
socket and SSH~\cite{p11kit} (verified).

In Crostini, users report that \code{pcscd} and \code{gpg} work with a shared
key (verified, two reports). FIDO keys for SSH do not work there, because
libfido2 needs \code{hidraw} (Finding~\ref{find:nohidraw}).

\subsection{Other approaches (E1 to E6)}

\paragraph{E1. Hybrid transport} Chrome can use a phone as authenticator.
The hybrid transport needs a Bluetooth advertisement near the
client~\cite{ctap22} (verified). A remote host has no phone nearby.

\paragraph{E2. CRD forwarding} Chromium has code to proxy security keys and
WebAuthn through CRD. Its policies are internal to Google, or for managed
users only, and the WebAuthn channel accepts only the official CRD extension
IDs~\cite{crdpol} (verified).

\paragraph{E3. WebAuthn proxy extension} The extension API
\code{chrome.webAuthenticationProxy} lets an extension take over WebAuthn in
Chrome~\cite{wap} (verified). auth-relay~\cite{authrelay} uses it on the
remote host and sends CTAP2 requests over SSH to python-fido2 on the laptop.
The proxied request has no origin field, so auth-relay guesses the origin
from the active tab (verified, source). Its laptop side needs \code{hidraw}.

\paragraph{E4. Smart card relay} vsmartcard~\cite{vsmartcard} (GPL-3.0)
relays a local PC/SC card to a virtual reader on another host (verified). It
can carry PIV and OpenPGP (assumption: untested with a YubiKey). The remote
host can then send any APDU to the card.

\paragraph{E5. Log in on the Chromebook} The Chromebook's own Chrome can open
a web app that the remote host serves, through an SSH SOCKS proxy. The host
name stays the same, so the relying party ID matches~\cite{webauthn3}
(assumption for the proxy setup). This avoids forwarding, but it does not help
when the login must happen inside the remote desktop.

\paragraph{E6. \code{vmc key-attach}} It shares the key with ChromeOS and
Linux at once, but only for U2F~\cite{crosvmdev} (verified). Passkeys need
FIDO2.
